\chapter{Testing y Control de Calidad}

\section{Estrategia de Testing}

El proyecto implementa una estrategia de testing en m\'ultiples niveles,
priorizando pruebas unitarias y funcionales con infraestructura m\'{\i}nima:

\section{Pruebas Unitarias (Vitest)}

\subsection{Configuraci\'on}

\begin{lstlisting}[style=typescript, caption=vitest.config.ts]
export default defineConfig({
  test: {
    globals: true,
    environment: 'node',
    include: ['src/**/*.test.{ts,tsx}'],
  },
});
\end{lstlisting}

\subsection{Pruebas Existentes}

\texttt{utils.test.ts} --- Pruebas para la funci\'on \texttt{cn()}:

\begin{lstlisting}[style=typescript, caption=Pruebas unitarias para cn()]
describe('cn', () => {
  it('combines class names', () => {
    expect(cn('foo', 'bar')).toBe('foo bar');
  });

  it('handles conditional classes', () => {
    expect(cn('base', false && 'hidden')).toBe('base');
  });

  it('merges tailwind classes correctly', () => {
    expect(cn('px-4', 'px-2')).toBe('px-2');
  });
});
\end{lstlisting}

\texttt{env.test.ts} --- Pruebas de variables de entorno:

\begin{lstlisting}[style=typescript, caption=Pruebas de variables de entorno]
describe('Environment Variables', () => {
  it('DATABASE_URL is defined', () => {
    expect(process.env.DATABASE_URL).toBeDefined();
  });
  it('AUTH_SECRET is defined', () => {
    expect(process.env.AUTH_SECRET).toBeDefined();
  });
});
\end{lstlisting}

\section{Pruebas Funcionales}

\subsection{Script de Testing via Bash}

El archivo \texttt{scripts/test.sh} ejecuta pruebas de API:

\begin{lstlisting}[style=bash, caption=test.sh - Pruebas funcionales de API]
#!/bin/bash
# Test de health endpoint
echo "Testing health endpoint..."
HEALTH=$(curl -s http://localhost:3000/api/health)
echo "$HEALTH" | grep -q '"status":"healthy"' && \
  echo "OK" || echo "FAIL"

# Test de rutas protegidas (deben redirigir a login)
echo "Testing protected routes..."
STATUS=$(curl -s -o /dev/null -w "%{http_code}" \
  http://localhost:3000/clients)
[ "$STATUS" = "302" ] && echo "OK" || echo "FAIL"

# Test de rutas publicas
echo "Testing public routes..."
STATUS=$(curl -s -o /dev/null -w "%{http_code}" \
  http://localhost:3000/auth/login)
[ "$STATUS" = "200" ] && echo "OK" || echo "FAIL"
\end{lstlisting}

\subsection{Script de Testing Integral}

El archivo \texttt{tests/comprehensive-test.ts} (366 l\'{\i}neas) realiza:

\begin{enumerate}
    \item \textbf{Fase 1:} Health check de la API.
    \item \textbf{Fase 2:} Verificaci\'on de endpoints sin autenticaci\'on.
    \item \textbf{Fase 3:} Autenticaci\'on y obtenci\'on de sesi\'on.
    \item \textbf{Fase 4:} Verificaci\'on de datos de seed en base de datos.
    \item \textbf{Fase 5:} Verificaci\'on de datos desde API autenticada.
    \item \textbf{Fase 6:} Conteos de registros en base de datos (40+
          aserciones).
\end{enumerate}

\subsection{Pruebas Manuales}

El archivo \texttt{tests/api.http} contiene 20+ ejemplos de peticiones
REST para VS Code REST Client, cubriendo todos los endpoints del sistema.

\section{CI/CD --- GitHub Actions}

\subsection{Workflow de CI}

\begin{lstlisting}[style=yaml, caption=.github/workflows/ci.yml]
name: CI
on:
  push:
    branches: [main]
  pull_request:
    branches: [main]

jobs:
  lint-typecheck:
    runs-on: ubuntu-latest
    steps:
      - uses: actions/checkout@v4
      - uses: pnpm/action-setup@v4
        with:
          version: 11
      - uses: actions/setup-node@v4
        with:
          node-version: 22
          cache: 'pnpm'
      - run: pnpm install --frozen-lockfile
      - run: pnpm lint
      - run: pnpm typecheck

  build:
    needs: lint-typecheck
    runs-on: ubuntu-latest
    steps:
      - uses: actions/checkout@v4
      - uses: pnpm/action-setup@v4
        with:
          version: 11
      - uses: actions/setup-node@v4
        with:
          node-version: 22
          cache: 'pnpm'
      - run: pnpm install --frozen-lockfile
      - run: pnpm build
\end{lstlisting}

\section{Pipeline Propuesto (Futuro)}

El plan de arquitectura describe el pipeline completo:

\begin{verbatim}
Push --> Lint + TypeCheck --> Test (Vitest)
    --> Build --> Deploy Staging --> E2E Tests (Playwright)
    --> Deploy Production
\end{verbatim}

\begin{itemize}[leftmargin=*]
    \item \textbf{E2E Tests:} Playwright para pruebas de interfaz gr\'afica.
    \item \textbf{Load Testing:} k6 o Artillery para pruebas de carga.
    \item \textbf{Deploy Staging:} Despliegue autom\'atico en entorno de
          pruebas.
    \item \textbf{Production Deploy:} Aprobaci\'on manual + deploy a Vercel
          o Docker.
\end{itemize}
